% Part: model-theory
% Chapter: models-of-arithmetic
% Section: standard-models

\documentclass[../../../include/open-logic-section]{subfiles}

\begin{document}

\olfileid{mod}{mar}{stm}
\olsection{算术的标准模型}

算术语言~$\Lang{L_A}$ 显然是为了讨论数，特别是自然数而设计的。因此，“那个”标准模型~$\Struct{N}$ 是特殊的：它是我们想要讨论的模型。但在逻辑中，我们通常只对结构性质感兴趣，而任何两个同构的结构都共享这些性质。所以，我们可以稍微放宽标准，把任何与~$\Struct{N}$ 同构的结构都视为“标准的”。

\begin{defn}
一个 $\Lang{L_A}$ 结构称为\emph{标准的}，当且仅当它与~$\Struct{N}$ 同构。
\end{defn}

\begin{prop}
\ollabel{prop:standard-domain} 若结构~$\Struct{M}$ 是标准的，则其论域是标准数码的值的集合，即，
\[
\Domain{M} = \Setabs{\Value{\num{n}}{M}}{n \in \Nat}
\]
\end{prop}

\begin{proof}
显然，每个 $\Value{\num{n}}{M} \in \Domain{M}$。我们只需证明每个 $x \in \Domain{M}$ 都等于某个 $n$ 对应的 $\Value{\num{n}}{M}$。
由于 $\Struct{M}$ 是标准的，它同构于~$\Struct{N}$。假设 $g\colon \Nat \to \Domain{M}$ 是一个同构映射。那么 $g(n) = g(\Value{\num{n}}{N}) =
\Value{\num{n}}{M}$。但由于 $g$ 是满射，对于每个 $x \in \Domain{M}$，都存在某个 $n \in \Nat$ 使得 $g(n) = x$。
\end{proof}

\begin{explain}
如果 $\Lang{L_A}$ 的一个结构~$\Struct{M}$ 是标准的，其论域中的所有元素都可以用标准数码 $\num{0}$、$\num{1}$、$\num{2}$……来命名，也就是那些项 $\Obj{0}$、$\Obj{0}'$、$\Obj{0}''$ 等等。当然，这并不意味着 $\Domain{M}$ 的元素\emph{就是}那些数，而只是说我们可以用和从~$\Domain{N}$ 中挑出数相同的方式来挑出它们。
\end{explain}

\begin{prob}
证明 \olref[mod][mar][stm]{prop:standard-domain} 的逆命题不成立，即给出一个结构~$\Struct{M}$ 的例子，使得 $\Domain{M} = \Setabs{\Value{\num{n}}{M}}{n \in \Nat}$，但它不同构于~$\Struct{N}$。
\end{prob}

\begin{prop}
\ollabel{prop:thq-standard}
如果 $\Sat{M}{\Th{Q}}$，并且 $\Domain{M} = \Setabs{\Value{\num{n}}{M}}{n
  \in \Nat}$，那么 $\Struct{M}$ 是标准的。
\end{prop}

\begin{proof}
我们需要证明 $\Struct{M}$ 同构于~$\Struct{N}$。考虑由 $g(n) = \Value{\num{n}}{M}$ 定义的函数 $g\colon \Nat \to \Domain{M}$。根据假设，$g$ 是满射。它也是单射：只要 $n \neq m$，就有 $\Th{Q} \Proves
\eq/[\num{n}][\num{m}]$。因此，由于 $\Sat{M}{\Th{Q}}$，只要 $n \neq m$，就有 $\Sat{M}{\eq/[\num{n}][\num{m}]}$。于是，若 $n \neq m$，则 $\Value{\num{n}}{M} \neq
\Value{\num{m}}{M}$，即 $g(n) \neq g(m)$。

我们还需要验证 $g$ 是一个同构映射。
\begin{enumerate}
\item 我们有 $g(\Assign{\Obj{0}}{N}) = g(0)$，因为 $\Assign{\Obj{0}}{N} = 0$。根据~$g$ 的定义，$g(0) =
  \Value{\num{0}}{M}$。但 $\num{0}$ 就是 $\Obj{0}$，而恰好为常项符号的项，其值由结构赋予该常项符号的对象给出，即 $\Value{\Obj{0}}{M} = \Assign{\Obj{0}}{M}$。所以，如所要求的那样，我们有 $g(\Assign{\Obj{0}}{N}) = \Assign{\Obj{0}}{M}$。
\item 在 $\Struct{N}$ 中，$\prime$ 是~$\Nat$ 上的后继函数，所以 $g(\Assign{\prime}{N}(n)) = g(n+1)$。接着，根据~$g$ 的定义，$g(n+1) =
  \Value{\num{n+1}}{M}$。但 $\num{n+1}$ 与 $\num{n}'$ 是同一个项，所以 $\Value{\num{n+1}}{M} =
  \Value{\num{n}'}{M}$。根据值函数的定义，这就是 $= \Assign{\prime}{M}(\Value{\num{n}}{M})$。由于 $\Value{\num{n}}{M} = g(n)$，我们得到 $g(\Assign{\prime}{N}(n)) =
  \Assign{\prime}{M}(g(n))$。
\item 在 $\Struct{N}$ 中，$+$ 是~$\Nat$ 上的加法函数，所以 $g(\Assign{+}{N}(n,m)) = g(n+m)$。接着，根据~$g$ 的定义，$g(n+m) =
  \Value{\num{n+m}}{M}$。但 $\Th{Q} \Proves
  \num{n+m} = (\num{n} + \num{m})$，所以 $\Value{\num{n+m}}{M} =
  \Value{\num{n}+\num{m}}{M}$。根据值函数的定义，这就是 $=
  \Assign{+}{M}(\Value{\num{n}}{M},\Value{\num{m}}{M})$。由于 $\Value{\num{n}}{M} = g(n)$ 以及 $\Value{\num{m}}{M} = g(m)$，我们得到 $g(\Assign{+}{N}(n, m)) = \Assign{+}{M}(g(n), g(m))$。
\item $g(\Assign{\times}{N}(n, m)) = \Assign{\times}{M}(g(n), g(m))$：练习。
\item $\tuple{n,m} \in \Assign{<}{N}$ 当且仅当 $n < m$。若 $n < m$，则 $\Th{Q} \Proves \num{n} < \num{m}$，并且也有 $\Sat{M}{\num{n} <
  \num{m}}$。因此 $\tuple{\Value{\num{n}}{M}, \Value{\num{m}}{M}} \in
  \Assign{<}{M}$，即 $\tuple{g(n), g(m)} \in \Assign{<}{M}$。如果 $n \not< m$，则 $\Th{Q} \Proves \lnot \num{n} < \num{m}$，从而 $\Sat/{M}{\num{n} < \num{m}}$。因此，和之前一样，$\tuple{g(n), g(m)} \notin \Assign{<}{M}$。综合起来，我们得到：$\tuple{n,m} \in \Assign{<}{N}$ 当且仅当 $\tuple{g(n), g(m)} \in
  \Assign{<}{M}$。
\end{enumerate}
\end{proof}

\begin{explain}
函数~$g$ 是定义从 $\Nat$ 到任何其他 $\Lang{L_A}$-结构~$\Struct{M}$ 论域的映射的最自然的方式，因为每个这样的 $\Struct{M}$ 都包含由 $\num{0}$、$\num{1}$、$\num{2}$ 等命名的元素。所以并不奇怪，如果 $\Struct{M}$ 至少使得关于 $\num{n}$ 的某些基本陈述以与~$\Struct{N}$ 相同的方式为真，并且 $g$ 又是双射，那么 $g$ 就会成为一个同构映射。事实上，如果 $\Domain{M}$ 除了 $\num{n}$ 命名的那些元素外不包含其他元素，那么 $g$ 就是唯一的同构映射。
\end{explain}

\begin{prop}
\ollabel{prop:thq-unique-iso} 如果 $\Struct{M}$ 是标准的，那么 \olref{prop:thq-standard} 证明中的 $g$ 是从 $\Struct{N}$ 到~$\Struct{M}$ 的唯一同构映射。
\end{prop}

\begin{proof}
假设 $h\colon \Nat \to \Domain{M}$ 是 $\Struct{N}$ 与~$\Struct{M}$ 之间的一个同构映射。我们通过对~$n$ 归纳来证明 $g = h$。若 $n = 0$，则根据~$g$ 的定义有 $g(0) = \Assign{\Obj{0}}{M}$。但由于 $h$ 是同构映射，$h(0) =
h(\Assign{\Obj{0}}{N}) =\Assign{\Obj{0}}{M}$，故而 $g(0) = h(0)$。

现在考虑 $n+1$ 的情形。我们有
\begin{align*}
  g(n+1) & = \Value{\num{n+1}}{M} \text{ 根据~$g$ 的定义}\\
  & = \Value{\num{n}'}{M} \text{ 因为 $\num{n+1}\ident \num{n}'$}\\
  & = \Assign{\prime}{M}(\Value{\num{n}}{M})
    \text{ 根据 $\Value{t'}{M}$ 的定义}\\
  & = \Assign{\prime}{M}(g(n))  \text{ 根据~$g$ 的定义}\\
  & = \Assign{\prime}{M}(h(n)) \text{ 由归纳假设}\\
  & = h(\Assign{\prime}{N}(n)) \text{ 由于 $h$ 是同构}\\
  & = h(n+1)
\end{align*}
\end{proof}

\begin{explain}
对任意可数无穷集合~$M$，在 $\Nat$ 与 $M$ 之间存在双射，因此每个这样的集合~$M$ 都可能成为某个标准模型~$\Struct{M}$ 的论域。事实上，一旦你选取一个对象 $z \in M$ 和一个合适的函数 $s$ 分别作为 $\Assign{\Obj{0}}{M}$ 和 $\Assign{\prime}{M}$，那么 $+$、$\times$ 和 $<$ 的解释就已确定。只有那些既是单射又是满射的函数~$s\colon M \to M \setminus \{z\}$ 才适合在标准模型中作为~$\Assign{\prime}{M}$。$s$ 的值域不能包含~$z$，否则 $\lforall[x][\eq/[\Obj 0][x']]$ 将为假。该语句在~$\Struct{N}$ 中为真，因此 $\Struct{M}$ 也必须使其为真。函数~$s$ 必须是单射，因为~$\Struct{N}$ 中的后继函数~$\Assign{\prime}{N}$ 是单射，且“$\Assign{\prime}{N}$ 是单射”这一事实可由一个在~$\Struct{N}$ 中为真的语句表达。它必须是满射，因为否则就会存在某个 $x \in M \setminus \{z\}$ 不在~$s$ 的值域中，也就是说，语句 $\lforall[x][(\eq[x][\Obj 0] \lor
\lexists[y][\eq[y'][x]])]$ 将在~$\Struct{M}$ 中为假——但它在~$\Struct{N}$ 中为真。
\end{explain}


\end{document}